\chapter{Discussion}
\label{sec:origins-discussion}
% ===================================================================

\section{Relation to Existing Work}

Assembly Theory~\cite{cronin2023assembly} has attracted both
enthusiasm and criticism.
The enthusiasm is for the biosignature criterion: a single
computable quantity that distinguishes living from non-living matter.
The criticism is precisely the two gaps identified in
Section~\ref{sec:assembly}: the lack of a spatial framework for
copy number and the undefined identity criterion for copies.
The present framework addresses both, though not in one step.
Namespace supplies the region and bisimulation supplies the criterion,
which makes the count well posed; but bisimulation is undecidable, so
the count is not yet one any instrument can take, and
Chapter~\ref{sec:scope-manufacture} has to relativise it to the
resolution an observer can afford before it becomes a measurement rather
than a definition.
The consequence is that copy number is observer-relative in a way
Assembly Theory does not anticipate, and that coarse instruments
overcount.

Smith's geosphere proposal~\cite{smith2016origin} is grounded in
thermodynamics and network autocatalysis rather than computation.
The connection we draw is not a reduction of Smith's proposal to
computation but a \emph{parallel structure}: the phase transition
in thermodynamic space corresponds to the basin-crossing in
process space, and the thermodynamic stability of the autocatalytic
network corresponds to the compactness of the coherence cluster.
Crucially, Smith's observation that metabolism is fundamentally
about the redistribution of energy from the solar gradient through
an ecosystem --- rather than merely replenishment of individual
agents --- is captured precisely by the account-flow graph of
Section~\ref{sec:metabolism}: producers couple to the gradient,
consumers redistribute the stored account, and the ecosystem is
viable when inflow exceeds dissipation.
The Prigogine notion of dissipative structures sustained far from
equilibrium maps directly onto the CPT-violating gradient-coupled
rules: a living agent is one that maintains a broken symmetry
between its forward and backward transition amplitudes by
continuously processing an external gradient.
Whether these are the same phenomenon viewed from two angles or
genuinely distinct mechanisms is a question for future work.

The connection between the Rho calculus Y combinator and
biological replication echoes earlier work on process calculi
as models of biological systems~\cite{regev2001representation},
but with a sharper mathematical claim: replication is not merely
\emph{modeled} by a fixed point but \emph{is} the fixed point
of the duplication map in any sufficiently expressive interactive
GSLT.

\section{The Deeper Implication}

The origins of life argument, if completed, would establish
something philosophically striking: that the emergence of life
is not a contingent accident of terrestrial chemistry but a
\emph{structural inevitability} of any sufficiently large
interactive computational system coupled to an external gradient.

Any universe described by a Turing-complete interactive GSLT
that is driven away from equilibrium by an external energy source
--- a gradient that breaks the CPT symmetry of the isolated
system --- will, given sufficient time and population, produce
metabolically viable replicating fixed points.
The density of replicators in the logical topology means there is
no avoiding them.
The gradient means there is always account available to sustain
them once they appear.
The phase transition to life is not an unlikely event; it is
the inevitable consequence of wandering in a space dense with
attractors, in the presence of a broken symmetry that keeps
the account replenished.

What emerges from the transition is not merely a replicating
process but an entire ecosystem: a compact coherence cluster
with a gradient-coupled producer layer, a food web of account
redistribution, and the beginnings of the trophic hierarchy
that is causal depth in the metabolic open synchronization tree.

This connects the whole book into a single narrative.
The machinery of Part~\ref{part:machinery} supplies the meter without
which none of the accounting below is well posed.
The learners of Part~\ref{part:ecology} are what the phase transition
produces, and the medium tower built there is what
Chapter~\ref{sec:assembly-depth} bounds by the assembly index, and what
Chapter~\ref{sec:scope-manufacture} shows must be populated at a cost
exponential in its height.
The physics of Part~\ref{part:physics} governs the wandering and
provides the decoration in which gradient coupling is a broken symmetry.
The cluster ontology of Part~\ref{part:sciencebot} provides the
coherence clusters that become ecosystems.
The choice and consciousness of Part~\ref{part:choice} describe what
those agents can come to know and feel, and how much of it their
assembly permits.
And the present part explains how they came to exist and how they sustain
themselves.

% ===================================================================
